\subsection{等变横截性}\label{sec:14c}
接下来要解决的问题，是通过上链层次的朴素群作用实现 $\Z/2$ 对称性。我们选择施加一系列几何约束，以简化这一问题的解决，而不追求最大的普遍性。

\StatementNumber{3}\begin{assumption}\label{st:14.3}
对于给定的带角正合辛流形 $M$ 及其对合 $\iota_M$，要求不动点集 $M^{\mathrm{fix}}$ 的实余维为二，且其法丛 $\nu$ 平凡。实际上，还要固定该丛的一个复平凡化 $\nu\to\CC$，并加上下述部分可积性条件：这一平凡化应是某个映射 $\pi:U\to\CC$ 的导数，其中 $U\subset M$ 是不动点集的某个邻域；映射满足 $\pi\circ\iota_M=-\pi$，且相对于 $I_M$ 计算的 $\bar\partial\pi$ 沿 $M^{\mathrm{fix}}$ 连同一阶导数一起消失。
\end{assumption}

假设给定 $I_M$ 及平凡化 $D\pi|M^{\mathrm{fix}}$。若 $N(I_M)$ 是 Nijenhuis 张量，则只要存在具有上述性质的 $\pi$，便有 $D\pi\circ N(I_M)$ 在 $M^{\mathrm{fix}}$ 上消失；这解释了“部分可积性”这一名称。\znote{203}{“部分可积性”指 $I_M$ 只在法向上可积。映射 $\pi$ 在 $M^{\mathrm{fix}}$ 上只被要求到一阶全纯（$\bar\partial\pi$ 及其一阶导数消失），因为 $M^{\mathrm{fix}}$ 是实子流形，而紧流形上整体全纯的 $\pi$ 只能是常值。存在这样的 $\pi$ 迫使法向部分的 Nijenhuis 张量消失，即 $D\pi\circ N(I_M)=0$；这正是复结构在 $M^{\mathrm{fix}}$ 的复法向方向上可积的表现。若 $I_M$ 沿法向也不可积，则不存在这样的 $\pi$；所以假设 14.3 实质上是对 $I_M$ 的一个限制，而非纯技术性的方便条件。[来源：原书 §(14c) 假设 14.3 及其后的说明，印刷页 203–204。]}

\StatementNumber{4}\begin{assumption}\label{st:14.4}
在所有 $\iota_M$ 不变的正合 Lagrangian 子流形 $L$ 中，只考虑完全不与不动点集相交，或从恒定方向趋近不动点集的那些。后一个条件是指，存在某个 $\sigma\in\R$，称为 $L$ 的\emph{法向角}，使对所有 $x\in L^{\mathrm{fix}}$ 都有
\[
D\pi(TL_x)=e^{i\sigma}\R\subset\CC.
\]
$\sigma$ 的可能取值有 $\Z$ 个，彼此相差 $\pi$ 的整数倍，但假设已经为每个有关的 Lagrangian 子流形选定其中一个。
\end{assumption}

继续讨论之前，需要引入一些术语。令 $H$ 为 $M$ 上的 $\iota_M$ 不变函数。在每一点 $x\in M^{\mathrm{fix}}$，法向的 Hessian $D^2H$
\SourcePage{212}
是 $\nu_x$ 上良定义的二次型，因为一阶导数必须消失。若对 $Z\in\nu_x$ 有
\begin{equation}\label{eq:14.9}
(D^2H)(Z)=h\,|Z|^2,
\end{equation}
则称 $H$ 以速度 $h\in\R$ \emph{无穷小地旋转} $\nu_x$。这里的 $|\cdot|$ 是将 $\nu$ 与 $TM|M^{\mathrm{fix}}$ 的反不变部分认同所得的 Hermite 度量。等价地，取 Hamilton 向量场 $X_H$，经 $D\pi$ 投影得到函数 $D\pi(X_H):U\to\CC$；则\eref{14.9}是说，在 $x$ 附近的一阶展开中有 $D\pi(X_H)=ih\pi$。

\clearpage
\StatementNumber{5}\begin{assumption}\label{st:14.5}
取满足\stref{假设}{14.4}的一对不变正合 Lagrangian 子流形 $(L_0,L_1)$，两者均与不动点集相交，法向角分别为 $\sigma_0,\sigma_1$。考虑 $\Z/2$ 不变的 Floer 数据 $(H,J)$ 时，施加下述附加要求。首先，与对 $I_M$ 的假设一样，要求相对于任意 $J_t$ 计算的导数 $\bar\partial\pi$ 沿 $M^{\mathrm{fix}}$ 连同一阶导数一起消失。其次，令 $\phi$ 为 $H$ 的“流”，则应存在满足
\begin{equation}\label{eq:14.10}
\sigma_1-\sigma_0-h\in(-\pi,0)
\end{equation}
的 $h\in\R$，以及开子集 $V\subset[0,1]\times M^{\mathrm{fix}}$，使其对所有 $t\in[0,1]$ 均包含 $\{t\}\times\phi^t(L_0^{\mathrm{fix}})$，并使 $H(t,\cdot)$ 对每个 $(t,x)\in V$ 都以速度 $h$ 无穷小地旋转 $\nu_x$。
\end{assumption}

这里只要求在足够大的开子集上控制法向 Hessian，因为按照 Floer 数据的定义，$H(t,\cdot)$ 必须在 $\partial M$ 附近为零，故不可能在整个 $M^{\mathrm{fix}}$ 上满足\eref{14.9}。关键在于，对 $x\in L_0^{\mathrm{fix}}$，映射 $D\phi_x^t:\nu_x\to\nu_{\phi^t(x)}$ 是乘以 $e^{iht}$ 的旋转。特别地，$\phi^1(L_0)$ 从方向 $e^{i(\sigma_0+h)}\R$ 趋近不动点集。仍由 Floer 数据的定义，该方向必须与 $L_1$ 横截；单是这一点就已蕴含 $\sigma_1-\sigma_0-h\in\R\setminus\pi\Z$，当然，$\phi^1(L_0)\cap L_1\cap M^{\mathrm{fix}}=\varnothing$ 的情形除外。附加要求\eref{14.10}选定 $\R\setminus\pi\Z$ 的一个“室”，即一个连通分支。这一条件看起来可能相当空泛，因为总能改变 $\sigma_0$ 或 $\sigma_1$ 来得到所需结果；但同时考虑多个 Floer 数据时，它就有了实际意义。

具有这些性质的函数仍十分丰富。特别地，可以任意指定 $H$ 在不动点集上，以及在该集合任意小邻域之外的值，但仍须在边界附近为零，且当然要 $\iota_M$ 不变。对殆复结构也一样，理由如下。必要时先缩小 $M^{\mathrm{fix}}$ 的邻域 $U$，使 $\pi|U$ 成为具有辛纤维的浸没。切丛 $TM|U$ 分解成 $TM^v=\ker(D\pi)$ 及其辛正交补 $TM^h$。后者由 $D\pi$ 平凡化，且在 $M^{\mathrm{fix}}$ 上的限制恰为 $\nu$。由于 $I_M$ 与平凡化 $D\pi|M^{\mathrm{fix}}$ 相容，沿不动点集，$I_M$ 的 $(hv)$ 分量为零，$(hh)$ 分量等于 $i$。对 $\bar\partial\pi$ 的附加要求
\SourcePage{213}
是说，这两个条件还应延伸到一阶导数；实际上，对 $(hh)$ 分量，这一点已由 $\iota_M$ 对称性推出。寻找具有同样性质的其他殆复结构时，显然可以自由指定其沿 $M^{\mathrm{fix}}$ 的 $(vv)$ 分量。

\StatementNumber{6}\begin{lemma}\label{st:14.6}
令 $L_0,L_1$ 为不变正合 Lagrangian 子流形，$(H,J)$ 为\stref{假设}{14.5}中的不变 Floer 数据。取 $u\in\M_Z(y_0,y_1)^*$，其像完全包含在 $M^{\mathrm{fix}}$ 中。则 $u$ 在 $M$ 中正则，当且仅当它在 $M^{\mathrm{fix}}$ 中正则。
\end{lemma}
\begin{proof}
线性化算子是 $\Z/2$ 等变的，因此分解为
\[
D_{Z,u}=D^-_{Z,u}\oplus D^+_{Z,u}.
\]
不变部分容易辨认。具体地，将 $u$ 视为 $M^{\mathrm{fix}}$ 中关于 Lagrangian 子流形 $(L_0^{\mathrm{fix}},L_1^{\mathrm{fix}})$ 及 Floer 数据 $(H|M^{\mathrm{fix}},J|M^{\mathrm{fix}})$ 的 Floer 方程解，并记为 $u^{\mathrm{fix}}$，便有
\begin{equation}\label{eq:14.11}
D_{Z,u^{\mathrm{fix}}}=D^+_{Z,u}.
\end{equation}
对于反不变部分，可用 $D\pi$ 将有关向量丛认同为 $E=Z\times\CC\to Z$，边界上的子丛 $F\subset E|\partial Z$ 则成为 $F_{s,0}=e^{i\sigma_0}\R$、$F_{s,1}=e^{i\sigma_1}\R$。沿 $M^{\mathrm{fix}}$，精确到一阶有 $D\pi\circ J_t=i\circ D\pi$。此外，当 $|s|\gg0$ 时，$x=u(s,t)$ 接近 $\phi^t(L_0)$，因而 $(t,x)\in V$；这意味着 $H(t,\cdot)$ 的 Hamilton 向量场 $X(t,\cdot)$ 在 $x$ 处精确到一阶满足 $D\pi(X(t,x))=ih\pi(x)$。因此，在条带状端上，
\begin{equation}\label{eq:14.12}
D^-_{Z,u}=\partial_s+i\partial_t+h.
\end{equation}
可以变形该算子，使\eref{14.12}处处成立，从而得到 $\operatorname{index}(D^-_{Z,u})=0$，例如见\stref{引理}{11.9}。另一方面，\stref{引理}{11.5}在 $\mu=0$、$|\Sigma^-|=1$ 时断言 $D^-_{Z,u}$ 单射，故它必可逆。结合\eref{14.11}便得到结论。
\end{proof}

有了这一结果，按照常规即可证明：\stref{假设}{14.5}所定义的这一相当特殊类别中的 Floer 数据一般是正则的。给定其中一个数据，先只考察不动点集 $M^{\mathrm{fix}}$，作微小扰动，使其在该处正则。由\stref{引理}{14.6}，这保证 $M$ 中所有不变 Floer 轨线正则。其余轨迹可用标准横截性论证的一种变体处理，见\bcite{85}命题 5.13。\znote{205}{把线性化算子按 $\Z/2$ 分解，是处理等变横截性的标准手法：等变算子与对合交换，故核与余核都带 $\Z/2$ 分解，整体满射当且仅当两部分各自满射。不变部分恰是 $M^{\mathrm{fix}}$ 中相应问题的算子\eref{14.11}；反不变部分因指标为零且单射而自动可逆。于是 $M$ 中的正则性完全由 $M^{\mathrm{fix}}$ 中的正则性决定，这正是引理 14.6 的实质，也说明为何只须在不动点集上作扰动。[来源：原书 §(14c) 引理 14.6 及其证明，印刷页 205。]}

令 $(L_0,\ldots,L_d)$，$d\geq2$，为满足\stref{假设}{14.4}的 $\iota_M$ 不变正合 Lagrangian 子流形，每一个都与不动点集相交，法向角分别为 $\sigma_0,\ldots,\sigma_d$。假设已按\stref{假设}{14.5}为子流形对 $(L_0,L_d)$ 及
\SourcePage{214}
$(L_{k-1},L_k)$，$1\leq k\leq d$，选定 Floer 数据。考虑带有 Lagrangian 标记 $(L_0,\ldots,L_d)$ 的 $(d+1)$ 标点圆盘万有族 $\mathcal S^{d+1}\to\RG^{d+1}$，并为其配备 $\iota_M$ 不变的扰动数据 $(K,J)$，与之前选定的 Floer 数据相容。

\StatementNumber{7}\begin{lemma}\label{st:14.7}
令 $(r,u)\in\M^{d+1}(y_0,\ldots,y_d)=\M_{\mathcal S^{d+1}}(y_0,\ldots,y_d)$ 为非齐次伪全纯多边形，且 $u$ 的像完全包含在 $M^{\mathrm{fix}}$ 中。则 $u$ 在 $M$ 中正则，当且仅当它在 $M^{\mathrm{fix}}$ 中正则。
\end{lemma}
\begin{proof}
如前，须证明算子 $D_{\mathcal S^{d+1},r,u}$ 的反不变部分可逆。第一步仍将这一部分写为 $D^-_{S,u}:W^{1,p}(S,E,F)\to L^p(S,E)$，其中 $S=\mathcal S^{d+1}_r$、$E=S\times\CC$，且 $F|C_k=e^{i\sigma_k}\R$。在每个条带状端上，该算子的行为由 Floer 数据规定，即如\eref{14.12}有 $\epsilon_{S,k}^*D^-_{S,u}=\partial_s+i\partial_t+h_k$，其中 $h_k$ 满足相应的条件\eref{14.10}。随后结合给出指标为零的\stref{命题}{11.13}与\stref{引理}{11.5}即可。
\end{proof}

与\stref{引理}{14.6}之后的讨论一样，这说明一般选取 $\iota_M$ 不变扰动数据 $(K,J)$，就能使所有模空间 $\M_{\mathcal S^{d+1}}$ 正则。事实上，这次论证稍简单，因为 $(d+1)$ 标点圆盘没有自同构。

\StatementNumber{8}\begin{note}\label{st:14.8}
如\stref{笔记}{11.6}所述，可以证明\stref{引理}{11.5}的一个版本，对有关算子在条带状端上的行为无需如此精细的控制。这又允许削弱这里施加的某些条件，将其换成纯拓扑的对应条件，只要它们恰好足以保证反不变线性化算子的指标为零即可。沿另一个方向，可以考虑不动点集具有更高余维的 $\Z/2$ 作用；但那时必须保留对反不变算子的显式控制，事实上还须将这种控制扩展到整个 Riemann 曲面，因为次数论证已不再适用。
\end{note}
